%%%%%%%%%%  scan idx 39  |  folio 26  %%%%%%%%%%
\SGAPrintPage{26}%
first $\lambda$-ring defined here, if the characteristic of $k$ is $0$, but not in the general case.

\sgaStableTarget{sga6.e00.ex.2-3}3) \emph{Determination of $\underline{K}(G)$ and of $\underline{K}_r(G)$ in certain important cases}: We suppose that the field $k$ is algebraically closed, that the group $G$ is algebraic, \emph{affine connected} and ``globally \emph{reductive}'' (the radical of $G$ is isomorphic to $k^{*s}$), and we consider only rational linear representations. In characteristic $0$, the complete reducibility of the rational linear representations of $G$ shows that $\underline{K}(G)=\underline{K}_r(G)$ is the free group generated by the classes of irreducible rational representations. Now, if $T$ is a maximal torus of $G$, we have $\underline{K}(T)=\underline{K}_r(T)=\Z(\That)$, the algebra of the group $\That$ of rational characters of $T$; we have $\That\simeq\Z^r$ if $r$ is the rank of $G$ (this is independent of the characteristic). The Weyl group $W$ operates on $T$, hence on $\underline{K}(T)$, in a way that depends only on the operations of $W$ on the free group $\That$. We have an evident homomorphism of $\lambda$-rings, deduced from the injection of $T$ into $G$:
\sgaStableTarget{sga6.e00.eq.1-13}\begin{equation}\tag{1.13}
\underline{K}_r(G)\longrightarrow \underline{K}_r(T)^W=\underline{K}(T)^W,
\end{equation}
where $\underline{K}(T)^W$ denotes the set of fixed points of $W$ in $\underline{K}(T)$. The theory of linear representations of $G$\SGAAuthorFootnote{*}{Cf.\ Séminaire Chevalley 1956/58 Groupes de Lie algébriques (École Normale Supérieure).} then proves the following theorem:

\sgaStableTarget{sga6.e00.thm.1-1}\textbf{Theorem 1.1.} \emph{The homomorphism} (1.13) \emph{is an isomorphism.}

We thus see that $\underline{K}_r(G)$, once the ``type'' of $G$ is known (characterized by $W$ operating on a group $\Z^r$), does not depend on the characteristic, and that it is a special $\lambda$-ring. We have a canonical basis of $\underline{K}_r(G)$ corresponding to the orbits of $W$ in $\That$; another basis is obtained as follows:
%%%%%%%%%%  scan idx 40  |  folio 27  %%%%%%%%%%
\SGAPrintPage{27}%
let $(\rho_i)_{1\leq i\leq r'}$ be the representation classes of $G$ corresponding to the \emph{fundamental} representations of $G/\rad G$, and let $(\sigma_i)_{1\leq i\leq r''}$ be a basis of the group of rational homomorphisms from $G$ to $k^*$; we thus have $r'+r''=r$ (the rank of $G$).

\sgaStableTarget{sga6.e00.cor.unnamed}\textbf{Corollary:} \emph{The ring $\underline{K}_r(G)$ is identified with the subring}
\[
\Z[\rho_1,\ldots,\rho_{r'};\sigma_1,\ldots,\sigma_{r''};\sigma_1^{-1},\ldots,\sigma_{r''}^{-1}]
\]
\emph{of the field of rational functions in the $\rho_i$ and the $\sigma_j$ with coefficients in $\Q$}\SGAAuthorFootnote{*}{The derived group $G'$ of $G$ must be assumed simply connected for this statement to be correct.}.

In particular, $\underline{K}_r(G)$ has no divisors of $0$; if $G=\prod_i\Gl(n_i,k)$, the ring $\underline{K}_r(G)$ is made explicit as follows: let $\rho_i$ be the representation of $G$ deduced from the identity representation of $\Gl(n_i,k)$; if we then consider the field of rational functions over $\Q$ in indeterminates $\lambda^j(\rho_i)$ ($1\leq j\leq n_i$), then $\underline{K}_r(G)$ is identified with the subring generated by these indeterminates and the inverses of the $\lambda^{n_i}(\rho_i)$. This justifies our assertion that, at least in characteristic $0$, the operations $\lambda^i$ allow all tensor operations on vector bundles to be reconstructed; this is still true over an algebraically closed field of arbitrary characteristic, provided reduced classes are taken.

\sgaStableTarget{sga6.e00.rem.unnamed}\textbf{Remarks:} a) In example 1), the ring $\underline{K}(G)$ is augmented, the augmentation being defined by the degree of the representations, and, by passage to the quotient, this augmentation defines an augmentation on $\underline{K}_r(G)$. Likewise, the $\lambda$-rings considered in example 2) are augmented when $X$ is $k$-irreducible; in $\underline{K}(\xi(X))$, the augmentation is defined by the \emph{rank of a vector bundle}, that is, the dimension of its fibres.

%%%%%%%%%%  scan idx 41  |  folio 28  %%%%%%%%%%
\SGAPrintPage{28}%
b)
When the field $k$ is not algebraically closed, it is plausible that $\underline{K}(G)$ is a special $\lambda$-ring in characteristic $0$, and likewise for $\underline{K}_r(G)$ in all characteristics.

c) In nonzero characteristic, theorem 1.1 was kindly supplied to me by Chevalley.

\sgaStableTarget{sga6.e00.sec.3.occ2}\subsection*{\S 3. The $\lambda$-ring defined by a graded ring}

Let $A$ be a commutative graded ring with unit; for simplicity we shall assume that the degrees of $A$ are positive and that $A^0=\Z$. We denote by $\Ahat$ the product of the $A^i$ (it is a commutative ring with unit), and by $\Ahp$ the kernel of the evident augmentation $\Ahat\to A^0=\Z$. Then $1+\Ahp$ is a subgroup of the multiplicative group of invertible elements of $\Ahat$. Consider the abelian group
\sgaStableTarget{sga6.e00.eq.1-14}\begin{equation}\tag{1.14}
\widetilde A=\Z\times(1+\Ahp),
\end{equation}
whose elements will be denoted $[n,x]$ with $n\in\Z$ and $x=1+\sum_{i\geq 1}x^i\in1+\Ahp$ ($x^i\in A^i$). We shall write this group additively, thus:
\sgaStableTarget{sga6.e00.eq.1-15}\begin{equation}\tag{1.15}
[n,x]+[n',x']=[n+n',xx'].
\end{equation}
The zero element of this group is $[0,1]$. We are going to introduce on it an augmented $\lambda$-ring structure, compatible with its additive structure, the augmentation being $\varepsilon:[n,x]\longrightarrow n$. The product $[m,x][n,y]$ is completely characterized by being written $[mn,x^n y^m(x*y)]$, where $x*y$ is an element of $1+\Ahp$ whose components are expressed as universal polynomials with integer coefficients in the $x^i$ and the $y^i$:
\sgaStableTarget{sga6.e00.eq.1-16}\begin{equation}\tag{1.16}
(x*y)^i=Q_i(x^1,\ldots,x^i;y^1,\ldots,y^i)\qquad (i\geq1)
\end{equation}
%%%%%%%%%%  scan idx 42  |  folio 29  %%%%%%%%%%
\SGAPrintPage{29}%
the
polynomial $Q_i$ being isobaric of weight $i$. Moreover, $x*y$ must be distributive with respect to the ``addition'' $xy$ and, for ``linear'' factors, must reduce to the following:
\sgaStableTarget{sga6.e00.eq.1-17}\begin{equation}\tag{1.17}
[1,1+x^1][1,1+y^1]=[1,1+x^1+y^1].
\end{equation}
The polynomials $Q_i$ are evaluated by calculations with elementary symmetric functions (this is the calculation of the Chern classes of a tensor product of vector bundles by means of the Chern classes $x^i$ and $y^i$ of the factors and their respective dimensions $m$ and $n$). From (1.17) we deduce:
\sgaStableTarget{sga6.e00.eq.1-17bis}\begin{equation}\tag{1.17 bis}
(1+x^1)*(1+y^1)=(1+x^1+y^1)/(1+x^1)(1+y^1),
\end{equation}
and an arbitrary product $x*y$ is calculated term by term by decomposing ``formally'' $x$ into a product $\prod_{i=1}^n(1+\alpha_i)$ and $y$ into a product $\prod_{i=1}^n(1+\beta_i)$, where the $\alpha_i$ and the $\beta_i$ have degree $1$.

In this way, $\widetilde A$ becomes an augmented commutative ring, with a unit $[1,1]$, and it can be interpreted as the ring obtained by adjoining a unit to the ring $1+\Ahp$. Moreover, on $\widetilde A$, we define a \emph{filtration} by putting $\widetilde A^0=\widetilde A$ and denoting by $\widetilde A^n$ the set of the $[0,x]$ with $x^i=0$ for $1\leq i<n$; this filtration is compatible with the product, for we have
\sgaStableTarget{sga6.e00.eq.1-18}\begin{equation}\tag{1.18}
\left(1+\sum_{i\geq m}x^i\right)*\left(1+\sum_{j\geq n}y^j\right)
=1-\frac{(m+n-1)!}{(m-1)!(n-1)!}x^m y^n+\cdots,
\end{equation}
the unwritten terms having degree $>m+n$. This formula shows that the graded ring $G(\widetilde A)$ associated with $\widetilde A$ is identified with $A$ from the additive point of view, but not from the multiplicative point of view; we can only say that the map $\eta'$ from $A$ to $G(\widetilde A)$ defined by
%%%%%%%%%%  scan idx 43  |  folio 30  %%%%%%%%%%
\SGAPrintPage{30}%
\sgaStableTarget{sga6.e00.eq.1-19}\begin{equation}\tag{1.19}
\eta'(x^i)=(-1)^{i-1}(i-1)!\,x^i\qquad (x^i\in A^i)
\end{equation}
is a \emph{ring homomorphism}.

We introduce on $\widetilde A$ operations $\lambda^i$ satisfying the axioms of $\lambda$-rings (\S 1, formulas (1.2)), well determined by the condition that the components of $\lambda^i[0,x]$ be deduced from the components $x^i$ of $x$ by \emph{universal} polynomials with integer coefficients
\sgaStableTarget{sga6.e00.eq.1-20}\begin{equation}\tag{1.20}
\lambda^i[0,x]=[0,\lambda^i x]\qquad (\lambda^i x)^{(n)}=Q_{i,n}(x^1,\ldots,x^n)
\end{equation}
($Q_{i,n}$ is isobaric of weight $n$), and that we have
\sgaStableTarget{sga6.e00.eq.1-21}\begin{equation}\tag{1.21}
\lambda^i[1,1+x^1]=0\qquad\text{for }i>1.
\end{equation}
The polynomials $Q_{i,n}$ are evaluated by a calculation with elementary symmetric functions, which is the calculation of the Chern classes of the exterior powers of a vector bundle, by means of the Chern classes and the rank of the latter. We then find that $\widetilde A$ is a \emph{special} augmented $\lambda$-ring and that the $\widetilde A^i$ are stable under the operations $\lambda^j$.

The augmented $\lambda$-ring $\widetilde A$ may be considered as a functor in $A$; consequently, the principal automorphism
\[
x=\sum_{i\geq0}x^i\longrightarrow \check x=\sum_{i\geq0}(-1)^i x^i
\]
of $A$ defines an automorphism of $\widetilde A$, denoted by the same symbol $\check{\ }$.

\sgaStableTarget{sga6.e00.prop.1-2}\textbf{Proposition 1.2.} \emph{For $x\in\widetilde A$ of the form $[n,1+x^1+\cdots+x^n]$} (we therefore have $x^i=0$ for $i>n$), \emph{we have $\lambda^i(x)=0$ for $i>n$, we have $\lambda^n(x)=[1,1+x^1]$, and finally}
\sgaStableTarget{sga6.e00.eq.1-22}\begin{equation}\tag{1.22}
\sum_{i=0}^n(-1)^i\lambda^i(x)=[0,1-(n-1)!x^n+\cdots]
\end{equation}
From this proposition and from the fact that the $\widetilde A^i$ are stable under the operations $\lambda^j$, we deduce that formula (1.22) is valid whenever $\varepsilon(x)=n$, whether or not we have $x^i=0$ for $i>n$.

